%=====================================================================
%  Preliminary pages
%=====================================================================
\begin{titlepage}
\centering
\vspace*{0.6cm}
{\Large\bfseries Integrating Multi-Source Data with Sentiment Analysis\\[4pt]
and Language Models to Enhance Stock Market\\[4pt] Decision Making}\\[1.6cm]
{\normalsize A THESIS}\\[0.4cm]
{\normalsize Submitted in partial fulfilment of the requirements\\
for the award of the degree of}\\[0.5cm]
{\large\bfseries DOCTOR OF PHILOSOPHY}\\[0.2cm]
{\normalsize in \Discipline}\\[1.4cm]
{\normalsize by}\\[0.35cm]
{\large\bfseries ANANDKUMAR PARDESHI}\\[0.18cm]
{\normalsize Registration No.\ \RegNo}\\[0.18cm]
{\normalsize ORCID \ScholarORCID}\\[1.4cm]
{\normalsize under the guidance of}\\[0.35cm]
{\large\bfseries Dr.\ SUJATA DESHMUKH}\\[0.18cm]
{\normalsize \GuideDesignation, Department of Computer Engineering}\\[0.18cm]
{\normalsize ORCID \GuideORCID}\\[1.6cm]
{\normalsize \ResearchCentre}\\[0.25cm]
{\normalsize \Univ}\\[1.6cm]
{\normalsize \SubMonth}
\vfill
{\footnotesize\color{phcolor} Items shown in red are placeholders; edit
\texttt{placeholders.tex} to replace them.}
\end{titlepage}

%---------------------------------------------------------------------
\chapter*{Declaration}
\addcontentsline{toc}{chapter}{Declaration}
\thispagestyle{plain}

I declare that this written submission represents my ideas in my own words and that,
where the ideas or words of others have been included, I have adequately cited and
referenced the original sources. I also declare that I have adhered to all principles
of academic honesty and integrity and have not misrepresented, fabricated or falsified
any idea, data, fact or source in my submission. I understand that any violation of
the above will be cause for disciplinary action by the University and may also evoke
penal action from the sources which have thus not been properly cited, or from whom
proper permission has not been taken when needed.

The work reported in this thesis has been carried out by me under the supervision of
Dr.~Sujata Deshmukh and has not been submitted, in whole or in part, for the award of
any other degree or diploma of this or any other institution.

\vspace{2.5cm}
\begin{flushright}
\begin{tabular}{@{}c@{}}
\rule{6cm}{0.4pt}\\
\textbf{Anandkumar Pardeshi}\\
Registration No.\ \RegNo\\
\SubMonth
\end{tabular}
\end{flushright}

%---------------------------------------------------------------------
\chapter*{Certificate}
\addcontentsline{toc}{chapter}{Certificate}
\thispagestyle{plain}

This is to certify that the thesis entitled \textbf{``Integrating Multi-Source Data with
Sentiment Analysis and Language Models to Enhance Stock Market Decision Making''},
submitted by \textbf{Anandkumar Pardeshi} (Registration No.\ \RegNo) to \Univ\ for the
award of the degree of \textbf{Doctor of Philosophy} in \Discipline, is a record of
bona fide research work carried out by him under my supervision and guidance at
\ResearchCentre.

The contents of this thesis, in full or in part, have not been submitted to any other
institute or university for the award of any degree or diploma. The thesis has reached
the standard of fulfilling the requirements of the regulations relating to the degree.

\vspace{2.5cm}
\begin{flushright}
\begin{tabular}{@{}c@{}}
\rule{6cm}{0.4pt}\\
\textbf{Dr.\ Sujata Deshmukh}\\
\GuideDesignation, Department of Computer Engineering\\
\ResearchCentre\\
\SubMonth
\end{tabular}
\end{flushright}

%---------------------------------------------------------------------
\chapter*{Acknowledgements}
\addcontentsline{toc}{chapter}{Acknowledgements}
\thispagestyle{plain}

I record my sincere gratitude to my research supervisor, Dr.~Sujata Deshmukh, whose
insistence on separating what the evidence supports from what it merely suggests shaped
this thesis more than any other single influence. Several of the results reported here
are negative, and they survive in this document because she treated an honest null as a
finding rather than a failure.

I thank the Research and Recognition Committee and the faculty of the Department of
Computer Engineering at \ResearchCentre\ for their assessment of this work at each
annual review, and my colleagues at the Department of Computer Science and Engineering,
Fr.~C.~Rodrigues Institute of Technology, Vashi, for accommodating the demands that
part-time doctoral research places on a teaching schedule.

I acknowledge the reviewers and editors of the venues in which the constituent studies
of this thesis were published or are under review. Their criticism materially improved
the work; in one instance a reviewer's insistence on a robustness analysis exposed two
implementation defects that had silently invalidated an entire results section, and
correcting them changed the central claim of that study.

Finally, I thank my family for their patience across the years this research has taken.

\vspace{1.5cm}
\begin{flushright}
\textbf{Anandkumar Pardeshi}
\end{flushright}

%---------------------------------------------------------------------
\chapter*{Abstract}
\addcontentsline{toc}{chapter}{Abstract}
\thispagestyle{plain}

Equity markets are driven simultaneously by structured price and volume records, by
qualitative information carried in news and social media, and by slower macroeconomic
conditions. Each stream is individually incomplete, and each becomes informative or
uninformative at different times. Conventional forecasting systems compress this
heterogeneity into a single numerical pipeline, report a statistical accuracy figure, and
stop short of the decision the investor actually has to take. This thesis addresses that
shortfall by developing, implementing and validating an integrated framework in which
multi-source fusion, language-model-based sentiment extraction, transparent reasoning and
decision-level allocation are treated as one connected engineering problem rather than as
four separate ones.

The first objective develops a comprehensive model integrating financial metrics, news
sentiment, social media sentiment and macroeconomic indicators. A hierarchical
Multi-Source Fusion Network is designed in which convolutional encoders process indicator
streams, transformer encoders process text, engineered temporal features carry
macroeconomic conditions, and a cross-modal attention layer assigns time-varying weight to
each source before a bidirectional recurrent encoder produces the forecast. Evaluated on
ten years of National Stock Exchange data across ten sectoral indices, the fused model
attains a normalised mean absolute error of 0.018, a coefficient of determination of 0.988
and a directional accuracy of 84.7 per cent at the index level, against 76.3 per cent for
a simple averaging ensemble over identical inputs. Removing both textual sources costs
10.9 per cent of explained variance, exceeding the sum of their individual removals and
confirming that sentiment carries information that price history does not. The fusion
mechanism is then strengthened by a Bayesian formulation in which the weight assigned to a
source is derived from its own recent predictive uncertainty through a softmax over
negative error variances, so that a stream is discounted at the moment it becomes
unreliable rather than at the end of a training cycle. On the Nifty~50 this dynamic rule
improves the Sharpe ratio by 22.7 per cent and reduces maximum drawdown by 18.3 per cent
against the strongest static-fusion baseline, while the same expert networks fused with
fixed weights perform no better than that baseline, which locates the gain in the weighting
mechanism rather than in the networks.

The second objective designs a hybrid framework coupling large language models with
conventional financial indicators. A transformer sentiment encoder converts news headlines
into a continuous polarity score, which is time-aligned to trading sessions and fused with
fifteen technical indicators before gradient-boosted classification. On Indian equity data
the sentiment score emerges as the highest-gain predictor, correlating 0.74 with momentum
yet only 0.04 with price, and the sentiment-augmented model raises directional accuracy
from a 46.43 per cent baseline average to 47.72 per cent while producing the highest
prediction confidence of all models examined. The absolute level is reported without
embellishment: at single-name daily horizons all models cluster near the random classifier,
so language-model sentiment supplies a measurable but bounded increment rather than a
decisive edge.

The third objective establishes transparency and examines sentiment noise and
sentiment-to-price temporal dynamics. A Dynamic-Causal Hybrid Framework replaces
correlation-based feature selection with a directed causal adjacency matrix recovered by
conditional independence testing, encodes intraday, daily and weekly horizons through a
multi-scale memory encoder, and allocates capital through a credibilistic optimiser. On
308 trading sessions of NIFTY50 data it recovers a causal graph of density 0.850 with an
identifiable causal hub, and demonstrates that strong causal influence can coexist with
weak linear correlation --- the mechanistic justification for preferring causal to
correlational feature selection. Sentiment quality is then treated as a measurable property
rather than an assumption: a sentiment composite is validated against four independently
published series, carries the expected sign against every one of them, and separates the
identified market states at $p = 1.8 \times 10^{-8}$.

The fourth objective evaluates and validates the framework as a decision system. Latent
market states are identified by an expanding-window Gaussian mixture model, the fused
signal enters a constrained quadratic program with covariance shrinkage and a turnover
penalty, and the resulting portfolio is scaled by a regime-dependent risk budget. Over
eighteen years of ten liquid multi-asset instruments under a walk-forward protocol with
transaction costs, the framework attains a Sharpe ratio of 0.878 against 0.628 for an
equal-weighted benchmark and reduces maximum drawdown from $-36.1$ to $-19.0$ per cent. It
outperforms in all four stress episodes, wins on drawdown in eighteen of nineteen calendar
years and in all eighteen specifications examined, and on a 665-session holdout that
postdates every design decision it records a higher Sharpe ratio and a smaller drawdown
than all three benchmarks. Ablation at matched average exposure isolates the source of the
gain: timing risk by market state improves both the Sharpe ratio and the drawdown with no
change in mean return, whereas the signal fusion layer contributes nothing detectable to
cross-sectional selection. Cross-market transfer to a concentrated mega-capitalisation
universe reproduces the defensive behaviour but loses on risk-adjusted return, which bounds
applicability to universes with genuine dispersion in asset volatility.

Taken together, the six publications underlying this thesis demonstrate a coherent
progression from information integration, through language-model-based signal extraction
and transparent reasoning, to validated decision making. The component-level nulls are
retained rather than removed, because they identify precisely which layers of a fusion
architecture earn their place and which do not.

\vspace{0.6em}
\noindent\textbf{Keywords:} multi-source data fusion; sentiment analysis; large language
models; explainable artificial intelligence; causal discovery; regime detection; portfolio
construction; Indian equity market.

%---------------------------------------------------------------------
\clearpage
\renewcommand{\contentsname}{Table of Contents}
{\setstretch{1.0}\tableofcontents}

\clearpage
\renewcommand{\listfigurename}{List of Figures}
{\setstretch{1.0}\listoffigures}

\clearpage
\renewcommand{\listtablename}{List of Tables}
{\setstretch{1.0}\listoftables}

%---------------------------------------------------------------------
\clearpage
\chapter*{List of Abbreviations}
\addcontentsline{toc}{chapter}{List of Abbreviations}
\thispagestyle{plain}

\begin{center}
\small
\setstretch{1.06}
\begin{tabular}{@{}p{2.1cm}p{4.6cm}p{2.1cm}p{4.6cm}@{}}
ADX & Average Directional Index & MAE & Mean Absolute Error \\
ANOVA & Analysis of Variance & MDD & Maximum Drawdown \\
ARIMA & Autoregressive Integrated Moving Average & ML & Machine Learning \\
ATR & Average True Range & MLP & Multi-Layer Perceptron \\
AUC & Area Under the Curve & MOSPI & Ministry of Statistics and Programme Implementation \\
BERT & Bidirectional Encoder Representations from Transformers & MSE & Mean Squared Error \\
Bi-LSTM & Bidirectional Long Short-Term Memory & MSFN & Multi-Source Fusion Network \\
BSE & Bombay Stock Exchange & NLP & Natural Language Processing \\
CAGR & Compound Annual Growth Rate & NSE & National Stock Exchange \\
CIT & Conditional Independence Test & OBV & On-Balance Volume \\
CNN & Convolutional Neural Network & OHLCV & Open, High, Low, Close, Volume \\
CPI & Consumer Price Index & PC & Peter--Clark (causal discovery algorithm) \\
CVaR & Conditional Value at Risk & QP & Quadratic Program \\
DA & Directional Accuracy & RBI & Reserve Bank of India \\
DCHF & Dynamic-Causal Hybrid Framework & RDM-DEA & Range Directional Measure Data Envelopment Analysis \\
DL & Deep Learning & RMSE & Root Mean Square Error \\
DM & Diebold--Mariano (test) & ROC & Receiver Operating Characteristic \\
EMH & Efficient Market Hypothesis & RSI & Relative Strength Index \\
ETF & Exchange-Traded Fund & SEBI & Securities and Exchange Board of India \\
GARCH & Generalised Autoregressive Conditional Heteroskedasticity & SHAP & Shapley Additive Explanations \\
GMM & Gaussian Mixture Model & SMA & Simple Moving Average \\
GRU & Gated Recurrent Unit & SVM & Support Vector Machine \\
IC & Information Coefficient & VADER & Valence Aware Dictionary and sEntiment Reasoner \\
IIP & Index of Industrial Production & VIX & Volatility Index \\
LLM & Large Language Model & VWAP & Volume-Weighted Average Price \\
LSTM & Long Short-Term Memory & WPI & Wholesale Price Index \\
MACD & Moving Average Convergence Divergence & XAI & Explainable Artificial Intelligence \\
\end{tabular}
\end{center}

%---------------------------------------------------------------------
\clearpage
\chapter*{List of Symbols}
\addcontentsline{toc}{chapter}{List of Symbols}
\thispagestyle{plain}

\begin{center}
\small
\setstretch{1.15}
\begin{tabular}{@{}p{3.0cm}L{10.0cm}@{}}
$h_m^{(l)}$ & representation of modality $m$ at encoder layer $l$ \\
$\alpha_m$ & cross-modal attention weight assigned to modality $m$ \\
$F_t$ & fused multi-source representation at session $t$ \\
$\sigma_{i,t}^2$ & predictive uncertainty of source $i$ at session $t$ \\
$w_{i,t}$ & uncertainty-derived fusion weight of source $i$ at session $t$ \\
$\tau$ & trailing window length used to estimate predictive uncertainty \\
$T$ & softmax temperature in the dynamic fusion rule \\
$s_t$ & continuous news-sentiment polarity score, $s_t \in [-1,+1]$ \\
$A(t)$ & time-varying directed causal adjacency matrix \\
$e_{ij}$ & directed causal edge from series $i$ to series $j$ \\
$r_t$ & identified latent market state (regime) at session $t$ \\
$T_{i,t},\, M_{i,t}$ & technical and sentiment composites for asset $i$ at session $t$ \\
$c_{i,t}$ & composite-agreement confidence term \\
$z_{i,t}$ & fused cross-sectional signal for asset $i$ \\
$\alpha_{i,t}$ & signal placed on a return scale by information-ratio scaling \\
$\Sigma_t$ & shrinkage-estimated return covariance matrix \\
$\lambda,\ \kappa$ & risk-aversion and turnover-penalty coefficients \\
$\tilde{w}_t,\ w_t^{*}$ & risky-sleeve weights before and after the regime risk budget \\
$k_{r_t}$ & regime risk-budget multiplier, $k \in \{1.00, 0.75, 0.50\}$ \\
$\tilde{R}$ & fuzzy (triangular) return variable in the credibilistic optimiser \\
\end{tabular}
\end{center}
